\subsection{Камера и видео в браузере}

Цель раздела: получить кадр с камеры дрона, показать его в браузере и нарисовать на нём телеметрию.

На дроне нет экрана, поэтому \verb|cv2.imshow| из главы~\ref{sec:python-vision} здесь не работает. Кадры берут из \verb|Camera|, а смотрят в браузере --- через \verb|ImageViewer|:

\begin{codefile}[python]{python-onboard/pioneer_sdk2/sdk2_06_camera_stream.py}
\end{codefile}

\paragraph{Что происходит.} \verb|Camera(camera_type=CameraType.MAIN)| открывает переднюю камеру; у Мини~2 есть ещё \verb|CameraType.OPT| --- нижняя камера оптического потока. \verb|get_cv_frame()| возвращает кадр --- массив numpy с цветами в порядке BGR, тот же формат, что в главе про камеру. \verb|viewer.imshow(name="stream", frame=frame)| отправляет кадр в видеопоток с именем \texttt{stream}; на компьютере, подключённом к Wi-Fi дрона, он открывается в браузере по адресу \texttt{http://10.42.0.1:8889/stream/} (видео идёт по WebRTC; если дрон подключён к вашей сети --- подставьте его адрес). Имя потока --- часть адреса, так что программа может показывать несколько потоков сразу. Путь кадра --- на рисунке~\ref{fig:sdk2-camera}.

\begin{tikzfig}
\begin{adjustbox}{max width=\linewidth}\begin{tikzpicture}[fig]
\node[fgrey, text width=29mm, minimum height=17mm] (cam) at (0,0) {\textbf{Камера дрона} \\ \texttt{Camera(...)}};
\node[forange, text width=29mm, minimum height=17mm] (fr) at (4.7,0) {\texttt{get\_cv\_frame()} \\ numpy, порядок BGR};
\node[forange, text width=29mm, minimum height=17mm] (cv) at (9.4,0) {Ваш код \\ \texttt{cv2}: рамки, надписи};
\node[forange, text width=29mm, minimum height=17mm] (vw) at (14.1,0) {\texttt{viewer.imshow(} \\ \texttt{name, frame)}};
\node[fblue, text width=29mm, minimum height=17mm] (cmd) at (9.4,-2.9) {Решение \\ команда дрону};
\node[fgreen, text width=29mm, minimum height=17mm] (br) at (14.1,-2.9) {Браузер: WebRTC; VLC: RTSP};
\draw[farrow] (cam) -- (fr);
\draw[farrow] (fr) -- (cv);
\draw[farrow] (cv) -- (vw);
\draw[farrow] (vw) -- (br);
\draw[fev] (cv) -- (cmd);
\end{tikzpicture}\end{adjustbox}
\captionof{figure}{Путь кадра на борту Мини~2. Экрана у дрона нет, поэтому кадр показывает браузер; тот же кадр можно использовать для решения (раздел о компьютерном зрении).}\label{fig:sdk2-camera}
\end{tikzfig}
 Тот же поток доступен по RTSP, например, для VLC: \texttt{rtsp://10.42.0.1:8554/stream}; адрес вида \texttt{rtsp://...:8889} не работает.

\begin{warnbox}[Останавливайте и поток, и камеру]
В конце программы нужно вызвать и \verb|viewer.close()|, и \verb|camera.stop()|. Если забыть \verb|camera.stop()|, поток приёма кадров продолжит работать, и программа после основного цикла не завершится --- такая ошибка была даже в примере из документации Геоскана.
\end{warnbox}

\begin{warnbox}[Окно пишет «stream not found»]
Это значит, что программа сейчас ничего не публикует: она ещё не запустилась, уже завершилась или имя потока в адресе не совпадает с первым аргументом \verb|imshow|. Страница сама переподключается каждые несколько секунд --- достаточно запустить программу.
\end{warnbox}

\begin{codebox}[Camera]
\verb|Camera(camera_type=CameraType.MAIN)| --- открыть камеру. Типы камер (\verb|CameraType.MAIN|, \verb|CameraType.OPT|) формируются из настроек платы: на плате без нижней камеры \verb|OPT| не будет.\\
\verb|camera.get_cv_frame(timeout=5.0)| --- следующий кадр (numpy, BGR); нет кадра за \verb|timeout| секунд --- обычно \verb|TimeoutError|, но на всякий случай проверяйте и \verb|None|.\\
\verb|camera.stop()| --- остановить захват. Ставьте в \verb|finally|.
\end{codebox}

Кадры копятся в очереди камеры, поэтому цикл с \verb|get_cv_frame()| идёт со скоростью камеры, а если ваша обработка медленнее, вы начинаете отставать от реальности. После резкого изменения сцены (например, наклона камеры) первые кадры могут оказаться старыми: пропустите несколько. На платах с драйвером gstreamer можно открыть несколько камер сразу (переднюю и нижнюю).

\begin{codebox}[ImageViewer]
\verb|ImageViewer()| --- создать показ.\\
\verb|viewer.imshow(name, frame, fps=30)| --- отправить кадр в поток \verb|name|; поток создаётся при первом кадре, и размер этого кадра задаёт размер потока, так что все следующие кадры того же имени должны быть того же размера. \verb|viewer.close()| --- закрыть все потоки. Ставьте в \verb|finally|.
\end{codebox}

\begin{infobox}[Камера на компьютере: драйвер rtsp]
Библиотека выбирает драйвер камеры по настройкам платы. На бортовых платах (Мини~2, Radxa Zero~3W) кадры идут из общей памяти по драйверу \texttt{gstreamer}, и \verb|ImageViewer| доступен. На платах с драйвером \texttt{rtsp} (так настроена конфигурация, которая поставляется с пакетом для Raspberry Pi Zero) \verb|Camera(camera_type, camera_ip="10.42.0.1:8554")| читает RTSP-поток по адресу вида \texttt{rtsp://10.42.0.1:8554/<имя камеры>}; адрес обязательно с портом, иначе \verb|ValueError|. \verb|ImageViewer| в этом режиме нет: показывать кадры придётся самому, например \verb|cv2.imshow| на компьютере.
\end{infobox}

\paragraph{Минимальная трансляция.} Самый короткий вариант --- 30~секунд берём кадры и отправляем в окно:

\begin{codefile}[python]{python-onboard/pioneer_sdk2/14a_camera_stream.py}
\end{codefile}

\paragraph{Рисуем поверх кадра (OSD).} Кадр --- обычный массив numpy, поэтому на нём можно рисовать средствами OpenCV: телеметрия, перекрестие, подписи.

\begin{codefile}[python]{python-onboard/pioneer_sdk2/14_camera_osd.py}
\end{codefile}

\paragraph{Что происходит.} Если кадр не пришёл вовремя, \verb|get_cv_frame| может выбросить \verb|TimeoutError| или вернуть \verb|None|, поэтому пример обрабатывает оба случая. Полупрозрачная подложка: прямоугольник рисуется на копии кадра и смешивается с оригиналом через \verb|cv2.addWeighted|. Частота кадров сглажена экспоненциальным средним: $\text{fps} \leftarrow 0.9\,\text{fps} + 0.1/\Delta t$. В \verb|finally| закрываются и трансляция, и камера.

\begin{taskbox}[Две камеры в одном окне]
Откройте обе камеры Мини~2 (\verb|CameraType.MAIN| и \verb|CameraType.OPT|) и транслируйте их в два потока \texttt{front} и \texttt{down}. На каждом кадре подпишите, откуда он (\verb|cv2.putText|). Закройте всё в \verb|finally|.
\expected{по адресам \texttt{.../front/} и \texttt{.../down/} (в симуляторе --- два окна) идёт видео с двух камер, подписи не путаются; после \texttt{Ctrl+C} программа останавливается без трассировки.}
\mission{Мини~2 SDK2}{30--32}
\end{taskbox}

\begin{warnbox}[Кадры не приходят]
Если \verb|get_cv_frame| всегда выбрасывает \verb|TimeoutError|, возможно, завис сервис захвата видео (в журнале ядра \texttt{dmesg} --- \texttt{not active buffer, skip frame}). Перезапустите его на дроне: \verb|sudo systemctl restart media-server|.
\end{warnbox}

\subsection{Наклон камеры: сервопривод}

Цель раздела: направлять камеру вниз или вперёд программой и получать корректные кадры после поворота.

Передняя камера Мини~2 стоит на сервоприводе и наклоняется вверх-вниз. Им управляет класс \verb|ServoCamera|:

\begin{codebox}[ServoCamera]
\verb|servo = ServoCamera()| --- подключиться к сервоприводу (SDK проверяет, что он есть на борту; на плате без него --- \verb|ValueError|).\\
\verb|servo.set_angle(angle, priority=ServoPriority.LOW)| --- наклонить камеру: \verb|angle| в градусах от $-80$ (почти вниз) до $30$ (вверх), $0$ --- прямо вперёд. Угол вне диапазона --- \verb|ValueError|. Возвращает \verb|True|, если команда отправлена; \verb|False|, если сервопривод не отвечает. \verb|priority| --- приоритет команды: \verb|ServoPriority.HIGH|, \verb|MEDIUM| или \verb|LOW|.
\end{codebox}

\begin{codefile}[python]{python-onboard/pioneer_sdk2/sdk2_07_servo_camera.py}
\end{codefile}

\paragraph{Что происходит.} Дрон взлетает и по очереди наклоняет камеру: вперёд, на $45^\circ$ вниз, почти отвесно вниз и обратно. После каждой команды программа две секунды показывает кадры с подписью угла --- видно, как в кадр входит площадка под дроном. \verb|set_angle| возвращается сразу, а сервопривод поворачивается за доли секунды, поэтому первые кадры после команды ещё «доезжают». Наклон меняет и кадры \verb|Camera(CameraType.MAIN)|: камеру, смотрящую вниз, можно использовать, чтобы искать метки на полу, не сажая дрон.

SDK только передаёт приоритет вместе с углом службе сервопривода на борту; как она разрешает спор между командами разных приоритетов, описание SDK не говорит. Если наклоном управляет одна программа, оставляйте \verb|LOW| по умолчанию. Если служба сервопривода недоступна (программа запущена не на дроне), \verb|set_angle| не бросает исключение, а возвращает \verb|False|.

\subsubsection*{Серия снимков под разными углами}

\begin{codefile}[python]{python-onboard/pioneer_sdk2/15_servo_panorama.py}
\end{codefile}

\paragraph{Что происходит.} Кадры копятся в очереди, поэтому сразу после поворота можно получить «старый» снимок --- \verb|fresh_frame()| пропускает несколько кадров. Миниатюры масштабируются с сохранением пропорций и склеиваются \verb|cv2.vconcat| в один «контактный лист»; формат \verb|f"{angle:+03d}"| даёт имена вида \texttt{tilt\_+30.jpg}, \texttt{tilt\_-50.jpg}. В \verb|finally| камера возвращается в горизонтальное положение. В симуляторе \verb|cv2.imwrite| --- заглушка: программа проходит до конца, но файлы не сохраняются.

\begin{taskbox}[Метки на полу]
Взлетите, наклоните камеру Мини~2 на $-80^\circ$ и показывайте поток через \verb|ImageViewer| с подписью текущего угла. Затем верните камеру в $0^\circ$ и сядьте. Если кадра нет 5~с, обработайте \verb|TimeoutError| и сядьте.
\expected{в окне камеры виден пол под дроном и подпись \texttt{angle -80}; после возврата --- снова вид вперёд; программа завершается посадкой.}
\mission{Мини~2 SDK2}{33}
\end{taskbox}

\begin{taskbox}[Плавный наклон]
Наклоняйте камеру от $0^\circ$ до $-80^\circ$ шагом $10^\circ$ каждые 0.5~с и обратно. Попробуйте угол $-90$ и обработайте исключение.
\expected{камера плавно опускается и поднимается; угол $-90$ вызывает \texttt{ValueError}, который программа ловит и печатает допустимый диапазон.}
\mission{Мини~2 SDK2}{33}
\end{taskbox}

\subsection{Запись видео и фото: RecorderControl}

Цель раздела: записывать видео полёта и снимать фото в заданных точках, не расходуя на это кадры вашей программы.

Класс \verb|RecorderControl| (Мини~2) пишет видео и снимает фото в галерею дрона. Запись идёт параллельно с полётом и не требует кадров в вашей программе.

\begin{codebox}[RecorderControl (по сборнику Геоскана, 0.16)]
\verb|recorder = RecorderControl()| --- подключиться к службе записи.\\
\verb|recorder.get_camera_recording_configs(camera_type)| --- список доступных режимов записи камеры (ширина, высота, fps).\\
\verb|recorder.start_recording(camera_type,| \verb|config, output_dir=None)| --- начать запись видео; по умолчанию файлы кладутся в \texttt{/mnt/media/videos/}.\\
\verb|recorder.stop_recording(camera_type)| --- закончить запись.\\
\verb|recorder.take_photo(file_name, camera_type,| \verb|config, output_dir=None)| --- сохранить снимок; расширение имени --- \texttt{.jpg}, \texttt{.jpeg} или \texttt{.png}. Пустое имя или неверное расширение --- \verb|ValueError|, ошибка камеры --- \verb|RuntimeError|.
\end{codebox}

\begin{codefile}[python]{python-onboard/pioneer_sdk2/16_recorder_mission.py}
\end{codefile}

\paragraph{Что происходит.} Программа печатает режимы записи основной камеры и берёт первый; \verb|config=None| означает, что библиотека выберет первый доступный. Видео пишется весь полёт, а в трёх вершинах после секундной паузы (чтобы погасить раскачку) делается фото. Флаг \verb|recording| гарантирует, что \verb|stop_recording()| вызывается, только если запись началась. Готовые файлы лежат в галерее медиа: \texttt{http://10.42.0.1:5000}, папки \texttt{photos} и \texttt{videos}.

\begin{infobox}[RecorderControl в симуляторе]
Симулятор проверяет аргументы, ведёт состояние записи (повторный старт --- \verb|RuntimeError|, \verb|stop_recording()| без записи --- \verb|False|) и печатает «сделал бы», но файлов не создаёт. Режимы основной камеры там: $1280\times720$ и $640\times360$ при 30 кадрах в секунду.
\end{infobox}

\begin{taskbox}[Облёт с видео]
Поставьте в центр площадки предмет, облетите его по окружности (пример 07) и запишите видео облёта через \verb|RecorderControl|. Фото сделайте в каждой четвёртой точке с подписью времени в имени файла.
\expected{в галерее дрона появилось видео всего облёта и фото с уникальными именами (в симуляторе --- соответствующие строки журнала); запись останавливается и при ошибке (в \texttt{finally}).}
\mission{Мини~2 SDK2}{34}
\end{taskbox}
